% ── chemistry-handout-preamble.tex ── 化学竞赛讲义 LaTeX 导言区（v7 · 精排版）
\documentclass[a4paper,11pt]{ctexart}

% ── 字体优化 ──
\setCJKmainfont[AutoFakeBold=true,BoldFont=SimHei]{SimSun}
\setCJKsansfont{SimHei}

% 化学式宏（mhchem）
\usepackage[version=4]{mhchem}

% ── 数学 ──
\usepackage{amsmath,amssymb}

% ── 图片 ──
\usepackage{graphicx}

% ── 表格（跨页+自动列宽）──
\usepackage{booktabs}
\usepackage{xltabular}
\usepackage{array}
\newcolumntype{L}{>{\raggedright\arraybackslash}X}
\newcolumntype{C}{>{\centering\arraybackslash}X}
\newcolumntype{R}{>{\raggedleft\arraybackslash}X}

% ── 浮动体（[H] 强制定位）──
\usepackage{float}

% ── 页面 ──
\usepackage[includeheadfoot,margin=2.2cm,headheight=14pt]{geometry}
\setlength{\textwidth}{16.6cm}

% ── 行间距 ──
\usepackage{setspace}
\setstretch{1.15}
\setlength{\parskip}{0.3em}
\setlength{\parindent}{2em}

% ── 防止表格跨页 ──
\usepackage{longtable}
\setlength{\LTpre}{0pt}
\setlength{\LTpost}{0pt}

% ── 防止图片跨页（尽量）──
\renewcommand{\textfraction}{0.1}
\renewcommand{\topfraction}{0.9}
\renewcommand{\bottomfraction}{0.9}
\setcounter{totalnumber}{10}
\setcounter{topnumber}{10}
\setcounter{bottomnumber}{10}

% ── 超链接(无红框) ──
\usepackage[colorlinks=true,linkcolor=black,citecolor=black,urlcolor=black]{hyperref}

% ── 颜色 ──
\usepackage[dvipsnames,svgnames]{xcolor}
\definecolor{chapterblue}{RGB}{23,50,77}
\definecolor{insightbg}{RGB}{235,240,250}
\definecolor{warningbg}{RGB}{255,245,235}
\definecolor{tipbg}{RGB}{235,250,235}

% ── 页眉页脚 ──
\usepackage{fancyhdr}
\pagestyle{fancy}
\fancyhf{}
\fancyhead[L]{\small\color{chapterblue}\leftmark}
\fancyhead[R]{\small\color{chapterblue}\thepage}
\renewcommand{\headrule}{{\color{chapterblue}\hrule height 0.8pt}}

% ── 章节编号与样式 ──
\setcounter{secnumdepth}{1}
\ctexset{
  section = {
    format = \Large\bfseries\heiti\color{chapterblue},
    aftername = \hspace{0.5em},
    beforeskip = 2ex,
    afterskip = 1ex,
  },
  subsection = {
    format = \large\bfseries\heiti,
    beforeskip = 1.5ex,
    afterskip = 0.5ex,
  },
  subsubsection = {
    format = \normalsize\bfseries\heiti,
    beforeskip = 1ex,
    afterskip = 0.3ex,
  },
  paragraph = {
    format = \small\bfseries,
    beforeskip = 0.5ex,
    afterskip = 0.2ex,
  },
}

% ── blockquote（引用/提示框）缩减 ──
\usepackage{etoolbox}
\AtBeginEnvironment{quote}{\small\setstretch{1.05}\leftskip=1em\rightskip=1em}

% ── 教学提示框（浅底色版）──
\newenvironment{insightbox}{%
  \par\smallskip
  \noindent{\small\bfseries\color{RoyalBlue}🧠 认知冲突}\par
  \leavevmode\colorbox{insightbg}{\parbox{\dimexpr\textwidth-2\fboxsep\relax}{%
  \ignorespaces}}%
}{\par\smallskip}

\newenvironment{warningbox}{%
  \par\smallskip
  \noindent{\small\bfseries\color{BrickRed}⚠️ 易错点}\par
  \leavevmode\colorbox{warningbg}{\parbox{\dimexpr\textwidth-2\fboxsep\relax}{%
  \ignorespaces}}%
}{\par\smallskip}

\newenvironment{tipbox}{%
  \par\smallskip
  \noindent{\small\bfseries\color{ForestGreen}💡 理解提示}\par
  \leavevmode\colorbox{tipbg}{\parbox{\dimexpr\textwidth-2\fboxsep\relax}{%
  \ignorespaces}}%
}{\par\smallskip}

% ── 列表（答案区加间距）──
\usepackage{enumitem}
\setlist{nosep}
% 专门用于参考答案的列表（每项之间空一行）
\newlist{answerlist}{enumerate}{1}
\setlist[answerlist]{label=\arabic*., leftmargin=2em, itemsep=0.8em}

% ── 防止内容溢出 ──
\usepackage{microtype}
\setlength{\emergencystretch}{2em}
\setlength{\hfuzz}{2pt}
\setlength{\overfullrule}{0pt}
\sloppy

% ── pandoc 兼容 ──
\providecommand{\tightlist}{\setlength{\itemsep}{0pt}\setlength{\parskip}{0pt}}
\newenvironment{Shaded}{}{}
\newenvironment{Highlighting}{}{}
\newcommand{\NormalTok}[1]{#1}
\newcommand{\CommentTok}[1]{#1}
\newcommand{\KeywordTok}[1]{#1}
\newcommand{\StringTok}[1]{#1}
\newcommand{\FunctionTok}[1]{#1}
\newcommand{\DataTypeTok}[1]{#1}
\newcommand{\NumberTok}[1]{#1}
\newcommand{\ControlFlowTok}[1]{#1}
\newcommand{\OperatorTok}[1]{#1}
\newcommand{\BaseNTok}[1]{#1}
\newcounter{none}

% ── 图片尺寸限制 ──
\AtBeginDocument{\setkeys{Gin}{width=\textwidth,keepaspectratio}}

% ── 防止图片跨节 ──
\usepackage{placeins}

% ── 自定义命令 ──
% 练习题标题（去掉括号）
% 例题标题（去掉中括号）
% 这些需要在convert.py预处理中做正则替换

\graphicspath{{C:/Obsidion/妙妙屋/11-模板/scripts/.chem_media/}}
\title{化学平衡}
\date{}
\setlength{\tabcolsep}{3.5pt}
\small
\begin{document}
\maketitle
\tableofcontents
\newpage
\section{化学平衡}\label{ux5316ux5b66ux5e73ux8861}

\textbf{本讲定位}：热力学告诉反应''能不能发生''，化学平衡进一步回答''能进行到什么程度''。平衡常数是连接热力学与实际化学体系的核心工具。
\textbf{上节衔接}：← 热力学初步-超级充实版(自学完整)\textbar 热力学初步(\(\Delta \1^{\circ}\) = −RT ln \(K^{\circ}\) 的热力学基础)
\textbf{下节衔接}：→ 2026-06-23-化学动力学基础\textbar 化学动力学(平衡到达的速率问题)→ 2026-06-23-酸碱理论\textbar 酸碱理论(弱酸弱碱平衡)→ 2026-06-23-电化学基础\textbar 电化学(\(\Delta \1^{\circ}\) = −nFE° 的电化学应用)
\textbf{主框架教材}：《普通化学原理》第4版 第6章
\textbf{辅助教材}：上海中学竞赛课程第四讲、Atkins《物理化学》化学平衡章节、质心化学原理L4

\begin{center}\rule{0.5\linewidth}{0.5pt}\end{center}

\subsection{学习目标}\label{ux5b66ux4e60ux76eeux6807}

\begin{itemize}
\tightlist
\item
  目标1：能正确书写各类反应的平衡常数表达式(\(K_{\1}\)、\(K_{\1}\)、\(K^{\circ}\))，掌握 \(K_{\1}\) = \(K_{\1}\)(RT)\^{}Δν 的推导与应用
\item
  目标2：能用 Q 与 K 的比较判断反应方向(Q \textless{} K 正向自发，Q \textgreater{} K 逆向自发)，理解 ΔG = RT ln(Q/\(K^{\circ}\)) 的推导
\item
  目标3：能用 ICE 表计算平衡转化率和平衡组成，处理气相与液相平衡问题
\item
  目标4：能运用勒夏特列原理和 Q-K 定量分析浓度、压力、温度变化对平衡的影响
\item
  目标5：能区分''改变 Q''(浓度/压力)与''改变 K''(温度)的本质差异，特别是恒容/恒压充惰性气体的辨析
\item
  目标6：能用 van't Hoff 方程定量计算温度对平衡常数的影响
\item
  目标7：能用 Clausius-Clapeyron 方程处理蒸气压与温度的关系、估算沸点
\end{itemize}

\begin{center}\rule{0.5\linewidth}{0.5pt}\end{center}

\subsection{一、为什么需要化学平衡}\label{ux4e00ux4e3aux4ec0ux4e48ux9700ux8981ux5316ux5b66ux5e73ux8861}

\textbf{动机段}：工业炼铁的主要反应为：
\[\mathrm{Fe_2O_3 + 3CO \rightleftharpoons 2Fe + 3CO_2}\]
按此方程式计算炼制 1 t 生铁需要多少焦炭，计算结果与实际情况有较大差别。因为在高炉中 CO 和 O\textsubscript{2} 不能全部转化，Fe\textsubscript{2}O\textsubscript{3} 和 CO 也不能全部转化为 Fe 和 CO\textsubscript{2}。也就是说，这些反应尽管可以自发发生(ΔG \textless{} 0)，但反应进行的程度是有限的。
化学平衡就是描述''反应能进行到什么程度''的理论框架。

\textbf{两类核心问题}：

{\def\LTcaptype{none} % do not increment counter
\begin{longtable}[]{@{}lll@{}}
\toprule\noalign{}
问题 & 回答者 & 本讲内容 \\
\midrule\noalign{}
\endhead
\bottomrule\noalign{}
\endlastfoot
反应能进行到什么程度？ & \textbf{化学平衡} & 平衡常数 K、转化率 α \\
反应进行的速率有多快？ & \textbf{化学动力学} & 速率方程、Arrhenius 方程 \\
\end{longtable}
}

\begin{center}\rule{0.5\linewidth}{0.5pt}\end{center}

\subsection{二、化学平衡的基本概念}\label{ux4e8cux5316ux5b66ux5e73ux8861ux7684ux57faux672cux6982ux5ff5}

\subsubsection{2.1 可逆反应与动态平衡}\label{ux53efux9006ux53cdux5e94ux4e0eux52a8ux6001ux5e73ux8861}

\textbf{可逆反应}：在相同条件下，正反应和逆反应都能进行的反应。注意可逆符号 ⇌ 与等号 = 的区别---可逆符号本身就暗示了平衡的存在。

\textbf{化学平衡的建立}：以 N\textsubscript{2}O\textsubscript{4} ⇌ 2NO\textsubscript{2} 体系为例。将 0.100 mol N\textsubscript{2}O\textsubscript{4} 充入 1 dm\textsuperscript{3} 密闭烧瓶，置于 373 K 恒温槽中。N\textsubscript{2}O\textsubscript{4}(无色)逐渐分解为 NO\textsubscript{2}(棕红色)，气体颜色逐渐变深。到一定时间后颜色不再加深，即达到平衡。取样分析得 {[}N\textsubscript{2}O\textsubscript{4}{]} = 0.040 mol/L，{[}NO\textsubscript{2}{]} = 0.120 mol/L。

从另一方向实验：将 0.100 mol NO\textsubscript{2} 充入同样烧瓶，颜色变浅(NO\textsubscript{2} 聚合为 N\textsubscript{2}O\textsubscript{4})，平衡后 {[}N\textsubscript{2}O\textsubscript{4}{]} = 0.014 mol/L，{[}NO\textsubscript{2}{]} = 0.072 mol/L。

两种起始条件不同，但平衡时 {[}NO\textsubscript{2}{]}\textsuperscript{2}/{[}N\textsubscript{2}O\textsubscript{4}{]} 的比值相同：

\[
K = \frac{[\mathrm{NO_2}]^2}{[\mathrm{N_2O_4}]} = 0.36 \quad (373\ \mathrm{K})
\]

{\def\LTcaptype{none} % do not increment counter
\begin{longtable}[]{@{}ll@{}}
\toprule\noalign{}
角度 & 平衡条件 \\
\midrule\noalign{}
\endhead
\bottomrule\noalign{}
\endlastfoot
\textbf{动力学} & 正反应速率 \(r_+\) = 逆反应速率 \(r_-\) \\
\textbf{热力学} & 恒温恒压、非体积功为零：\(\Delta_rG = 0\) \\
\textbf{宏观表现} & 各组分浓度(或分压)不再随时间改变 \\
\end{longtable}
}

\textbf{动态平衡的本质}：平衡时正逆反应仍在继续，只是速率相等。用同位素示踪可以验证---将 \textsuperscript{18}O 标记的 NO\textsubscript{2} 引入 N\textsubscript{2}O\textsubscript{4} ⇌ 2NO\textsubscript{2} 的平衡体系，最终 N\textsubscript{2}O\textsubscript{4} 中也会出现 \textsuperscript{18}O。这证明平衡不是反应''停止''，而是正逆反应''速度相等''。
类比理解：平衡就像一个进出人数相等的地铁站---人一直在进出，但站内总人数不变。

\subsubsection{2.2 平衡常数是反应的特征常数}\label{ux5e73ux8861ux5e38ux6570ux662fux53cdux5e94ux7684ux7279ux5f81ux5e38ux6570}

\textbf{平衡常数 K 与转化率 \(\alpha\) 的关键区别}：

{\def\LTcaptype{none} % do not increment counter
\begin{longtable}[]{@{}lccl@{}}
\toprule\noalign{}
量 & 与温度关系 & 与起始浓度关系 & 物理意义 \\
\midrule\noalign{}
\endhead
\bottomrule\noalign{}
\endlastfoot
\textbf{K} & 有关 & 无关 & 反应进行的限度(热力学本质) \\
\textbf{\(\alpha\)} & 有关 & 有关 & 特定条件下的转化程度 \\
\end{longtable}
}

以乙醇和醋酸酯化反应为例(373 K)：

{\def\LTcaptype{none} % do not increment counter
\begin{longtable}[]{@{}
  >{\centering\arraybackslash}p{(\linewidth - 8\tabcolsep) * \real{0.2000}}
  >{\centering\arraybackslash}p{(\linewidth - 8\tabcolsep) * \real{0.2000}}
  >{\centering\arraybackslash}p{(\linewidth - 8\tabcolsep) * \real{0.2000}}
  >{\centering\arraybackslash}p{(\linewidth - 8\tabcolsep) * \real{0.2000}}
  >{\centering\arraybackslash}p{(\linewidth - 8\tabcolsep) * \real{0.2000}}@{}}
\toprule\noalign{}
\begin{minipage}[b]{\linewidth}\centering
c(C\textsubscript{2}H\textsubscript{5}OH) / (mol/L)
\end{minipage} & \begin{minipage}[b]{\linewidth}\centering
c(CH\textsubscript{3}COOH) / (mol/L)
\end{minipage} & \begin{minipage}[b]{\linewidth}\centering
α(C\textsubscript{2}H\textsubscript{5}OH)
\end{minipage} & \begin{minipage}[b]{\linewidth}\centering
α(CH\textsubscript{3}COOH)
\end{minipage} & \begin{minipage}[b]{\linewidth}\centering
K
\end{minipage} \\
\midrule\noalign{}
\endhead
\bottomrule\noalign{}
\endlastfoot
3.0 & 3.0 & 67\% & 67\% & 4.0 \\
3.0 & 6.0 & 83\% & 42\% & 4.0 \\
6.0 & 3.0 & 42\% & 83\% & 4.0 \\
\end{longtable}
}

\textbf{结论}：K 在同一温度下恒定不变，但转化率随起始条件变化。当一种反应物过量时，\textbf{另一种反应物的转化率提高}，但\textbf{自身的转化率降低}。这是工业上''过量一种反应物''策略的热力学基础。

\begin{center}\rule{0.5\linewidth}{0.5pt}\end{center}

\subsection{三、平衡常数}\label{ux4e09ux5e73ux8861ux5e38ux6570}

\subsubsection{3.1 实验平衡常数}\label{ux5b9eux9a8cux5e73ux8861ux5e38ux6570}

对可逆反应 \(a\mathrm{A} + b\mathrm{B} \rightleftharpoons d\mathrm{D} + e\mathrm{E}\)，平衡时：

\[
$K_{\1}$ = \frac{[\mathrm{D}]^d[\mathrm{E}]^e}{[\mathrm{A}]^a[\mathrm{B}]^b} \qquad $K_{\1}$ = \frac{(p_\mathrm{D})^d(p_\mathrm{E})^e}{(p_\mathrm{A})^a(p_\mathrm{B})^b}
\]

\textbf{\(K_{\1}\) 与 \(K_{\1}\) 的关系推导}：

由理想气体方程 \(p_i = c_i RT = \frac{n_i}{V}RT\)，代入 \(K_p\) 表达式：

\[
$K_{\1}$ = \frac{([\mathrm{D}]RT)^d([\mathrm{E}]RT)^e}{([\mathrm{A}]RT)^a([\mathrm{B}]RT)^b} = \frac{[\mathrm{D}]^d[\mathrm{E}]^e}{[\mathrm{A}]^a[\mathrm{B}]^b} \cdot (RT)^{(d+e)-(a+b)}
\]

即：

\[
\boxed{$K_{\1}$ = $K_{\1}$(RT)^{\Delta\nu}}
\]

其中 \(\Delta\nu = (d+e) - (a+b)\)(气态生成物系数和 − 气态反应物系数和)。

\textbf{单位匹配}：\(R\) 的取值需与 \(p\)、\(c\) 单位匹配。\(p\) 用 kPa、\(c\) 用 mol/L 时，\(R = 8.314\ \mathrm{kPa \cdot L \cdot mol^{-1} \cdot K^{-1}}\)。若 \(p\) 用 bar，\(R = 0.0831\ \mathrm{bar \cdot L \cdot mol^{-1} \cdot K^{-1}}\)。

\subsubsection{3.2 书写平衡常数的规则}\label{ux4e66ux5199ux5e73ux8861ux5e38ux6570ux7684ux89c4ux5219}

\textbf{1. 纯固体和纯液体不写入表达式}：

\[
\mathrm{CaCO_3(s) \rightleftharpoons CaO(s) + CO_2(g)} \quad \Rightarrow \quad K = p_{\mathrm{CO_2}}
\]

固体浓度视为常数(活度 = 1)，已并入 K。又如：\(\mathrm{Fe_2O_3(s) + 3CO(g) \rightleftharpoons 2Fe(s) + 3CO_2(g)}\)，\(K_p = \frac{p^3(\mathrm{CO_2})}{p^3(\mathrm{CO})}\)，Fe\textsubscript{2}O\textsubscript{3} 和 Fe 都不出现在表达式中。

\textbf{2. 稀溶液中溶剂(水)不写入}：

\[
\mathrm{CH_3COOH(aq) + H_2O(l) \rightleftharpoons H_3O^+(aq) + CH_3COO^-(aq)}
\]

\[
$K_{\1}$ = \frac{[\mathrm{H_3O^+}][\mathrm{CH_3COO^-}]}{[\mathrm{CH_3COOH}]}
\]

原因：0.1 mol/L 的 HAc 溶液电离过程中消耗的 H\textsubscript{2}O 分子数占 H\textsubscript{2}O 总分子数的比例极小，{[}H\textsubscript{2}O{]} 可视为常数，归并入 K。
\textbf{非水溶剂例外}：酯化反应在乙醇中进行时，水的浓度不可视为常数。

\textbf{3. 同一反应不同写法 → 不同 K 值}：

{\def\LTcaptype{none} % do not increment counter
\begin{longtable}[]{@{}
  >{\raggedright\arraybackslash}p{(\linewidth - 4\tabcolsep) * \real{0.3077}}
  >{\raggedright\arraybackslash}p{(\linewidth - 4\tabcolsep) * \real{0.3077}}
  >{\centering\arraybackslash}p{(\linewidth - 4\tabcolsep) * \real{0.3846}}@{}}
\toprule\noalign{}
\begin{minipage}[b]{\linewidth}\raggedright
写法
\end{minipage} & \begin{minipage}[b]{\linewidth}\raggedright
K 表达式
\end{minipage} & \begin{minipage}[b]{\linewidth}\centering
373 K 数值
\end{minipage} \\
\midrule\noalign{}
\endhead
\bottomrule\noalign{}
\endlastfoot
\(\mathrm{N_2O_4 \rightleftharpoons 2NO_2}\) & \(K = [\mathrm{NO_2}]^2/[\mathrm{N_2O_4}]\) & 0.36 \\
\(\mathrm{\frac{1}{2}N_2O_4 \rightleftharpoons NO_2}\) & \(K' = [\mathrm{NO_2}]/[\mathrm{N_2O_4}]^{1/2}\) & 0.60 \\
\(\mathrm{2NO_2 \rightleftharpoons N_2O_4}\) & \(K'' = [\mathrm{N_2O_4}]/[\mathrm{NO_2}]^2\) & 2.8 \\
\end{longtable}
}

\textbf{规律}：方程式乘以系数 \(n\) → \(K' = K^n\)；方程式反向 → \(K' = 1/K\)。

\textbf{4. K 值必须注明温度}：

弱电解质电离常数(如 \(K_{\1}\)(HAc))在室温范围内随温度变化很小，可不注明。但气相反应的 K 随温度变化显著：

\[
K(\mathrm{N_2O_4 \rightleftharpoons 2NO_2}) = \begin{cases} 0.36 & (373\ \mathrm{K}) \\ 3.2 & (423\ \mathrm{K}) \end{cases}
\]

\subsubsection{\texorpdfstring{3.3 标准平衡常数 \(K^{\circ}\)}{3.3 标准平衡常数 K\^{}\{\textbackslash circ\}}}\label{ux6807ux51c6ux5e73ux8861ux5e38ux6570-kcirc}

用相对浓度 \(c/c^\theta\)(\(c^\theta = 1\) mol/L)或相对分压 \(p/p^\theta\)(\(p^\theta = 100\) kPa = 1 bar)表示：

\[
K^\theta = \frac{(c_\mathrm{D}/c^\theta)^d(c_\mathrm{E}/c^\theta)^e}{(c_\mathrm{A}/c^\theta)^a(c_\mathrm{B}/c^\theta)^b}
\]

\textbf{特点}：
- \textbf{无量纲}(纯数)，便于与热力学关联
- 与 Δ\_rG° 的关系：\(\Delta_rG^\circ = -RT\ln K^\theta\)
- 液相反应中 \(K^\theta\) 与 \(K_c\) \textbf{数值相等}(因 \(c^\theta = 1\))
- 气相反应中 \$K\^{}\theta = \(K_{\1}\) \cdot (p\textsuperscript{\theta)}\{-\Delta\nu\}\$

\textbf{实际工作中的简化}：实验平衡常数只要把实测分压除以 p° 用相对压力表示，实验平衡常数作为无量纲处理也是合理的。所以在实际工作中往往并不严格区分 K 和 \(K^{\circ}\)。

\subsubsection{\texorpdfstring{3.4 \(\Delta \1^{\circ}\) 与 \(K^{\circ}\) 的桥梁公式}{3.4 \textbackslash Delta \textbackslash1\^{}\{\textbackslash circ\} 与 K\^{}\{\textbackslash circ\} 的桥梁公式}}\label{delta-1circ-ux4e0e-kcirc-ux7684ux6865ux6881ux516cux5f0f}

\textbf{van't Hoff 等温式}(任意状态的 ΔG)：

对气相反应 \(m\mathrm{A(g)} + n\mathrm{B(g)} \rightleftharpoons q\mathrm{C(g)}\)：

\[
\Delta G(T) = \Delta G^\ominus(T) + RT\ln\frac{(p_\mathrm{C}/p^\ominus)^q}{(p_\mathrm{A}/p^\ominus)^m(p_\mathrm{B}/p^\ominus)^n}
\]

令反应商 \(Q = \frac{(p_\mathrm{C}/p^\ominus)^q}{(p_\mathrm{A}/p^\ominus)^m(p_\mathrm{B}/p^\ominus)^n}\)，则：

\[
\boxed{\Delta G(T) = \Delta G^\ominus(T) + RT\ln Q}
\]

\textbf{推导桥梁公式}：平衡时 ΔG = 0，Q = \(K^{\circ}\)，代入上式：

\[
0 = \Delta G^\ominus(T) + RT\ln K^\theta
\]

\[
\boxed{\Delta G^\ominus(T) = -RT\ln K^\theta \quad \Longleftrightarrow \quad \lg K^\theta = -\frac{\Delta G^\ominus(T)}{2.30RT}}
\]

\textbf{综合得方向判据}：

\[
\Delta G(T) = RT\ln\frac{Q}{K^\theta}
\]

{\def\LTcaptype{none} % do not increment counter
\begin{longtable}[]{@{}cccl@{}}
\toprule\noalign{}
条件 & Q/\(K^{\circ}\) & Δ\_rG & 反应方向 \\
\midrule\noalign{}
\endhead
\bottomrule\noalign{}
\endlastfoot
\(Q < K^\theta\) & \textless{} 1 & \textless{} 0 & \textbf{正向自发} \\
\(Q = K^\theta\) & = 1 & = 0 & \textbf{平衡} \\
\(Q > K^\theta\) & \textgreater{} 1 & \textgreater{} 0 & \textbf{逆向自发} \\
\end{longtable}
}

\textbf{实例：合成氨反应的 \(K^{\circ}\) 随温度变化}

查表得 298 K 时 Δ\_fG°(NH\textsubscript{3}, g) = −16.4 kJ/mol：

\[
\Delta_rG^\ominus(298\ \mathrm{K}) = 2 \times (-16.4) - 0 - 0 = -32.8\ \mathrm{kJ/mol}
\]

\[
\lg $K_{\1}$^\ominus = \frac{32800}{2.30 \times 8.314 \times 298} = 5.76 \quad \Rightarrow \quad $K_{\1}$^\ominus(298\ \mathrm{K}) = 5.8 \times 10^5
\]

在 673 K 时，利用 \(\Delta \1^{\circ}\) = \(\Delta \1^{\circ}\) − T\(\Delta \1^{\circ}\)(\(\Delta \1^{\circ}\) = −91.8 kJ/mol，\(\Delta \1^{\circ}\) = −0.198 kJ/(mol·K))：

\[
\Delta_rG^\ominus(673\ \mathrm{K}) = -91.8 - 673 \times (-0.198) = +41.5\ \mathrm{kJ/mol}
\]

\[
\lg $K_{\1}$^\ominus = \frac{-41500}{2.30 \times 8.314 \times 673} = -3.23 \quad \Rightarrow \quad $K_{\1}$^\ominus(673\ \mathrm{K}) = 5.9 \times 10^{-4}
\]

\textbf{工业启示}：298 K 时 \(K^{\circ}\) 很大(5.8 × 10\textsuperscript{5})，平衡时 NH\textsubscript{3} 分压远大于 N\textsubscript{2} 和 H\textsubscript{2}；但 673 K 时 \(K^{\circ}\) 很小(5.9 × 10\textsuperscript{-4})，平衡时 NH\textsubscript{3} 分压远小于 N\textsubscript{2} 和 H\textsubscript{2}。低温有利于平衡产率，但低温反应速率太慢。工业上选用 300\textsubscript{500°C(573}773 K)配铁催化剂，是平衡与速率的折中。
\textbf{\(\Delta \1^{\circ}\) 的粗估规则}：当 \(\Delta \1^{\circ}\) \textgreater{} +40 kJ/mol 时，\(K^{\circ}\) 极小，可认为反应不能正向进行；当 \(\Delta \1^{\circ}\) \textless{} −40 kJ/mol 时，\(K^{\circ}\) 极大，可认为反应能正向自发进行；\(\Delta \1^{\circ}\) 介于两者之间时，需结合具体条件分析。

\subsubsection{3.5 多重平衡规则}\label{ux591aux91cdux5e73ux8861ux89c4ux5219}

一种物质同时参与几种平衡，称为\textbf{多重平衡}。例如 SO\textsubscript{2}、SO\textsubscript{3}、NO、NO\textsubscript{2}、O\textsubscript{2} 共存时：

\begin{enumerate}
\def\labelenumi{\arabic{enumi}.}
\tightlist
\item
  \(\mathrm{SO_2 + \frac{1}{2}O_2 \rightleftharpoons SO_3}\)，\(K_{p_1}\)
\item
  \(\mathrm{NO_2 \rightleftharpoons NO + \frac{1}{2}O_2}\)，\(K_{p_2}\)
\item
  \(\mathrm{SO_2 + NO_2 \rightleftharpoons SO_3 + NO}\)，\(K_{p_3}\)
\end{enumerate}

由于反应 3. = 1. + 2.，而 \(\Delta \1^{\circ}\) 是广度量状态函数，所以：

\[
\Delta G_3^\ominus = \Delta G_1^\ominus + \Delta G_2^\ominus
\]

\[
\Rightarrow \quad \boxed{K_{p_3} = K_{p_1} \times K_{p_2}}
\]

\textbf{多重平衡规则总结}：

{\def\LTcaptype{none} % do not increment counter
\begin{longtable}[]{@{}ll@{}}
\toprule\noalign{}
反应式组合 & K 的运算 \\
\midrule\noalign{}
\endhead
\bottomrule\noalign{}
\endlastfoot
反应(3) = 反应(1) + 反应(2) & \(K_3 = K_1 \times K_2\) \\
反应(3) = 反应(1) − 反应(2) & \(K_3 = K_1 / K_2\) \\
反应(3) = n × 反应(1) & \(K_3 = K_1^n\) \\
\end{longtable}
}

\textbf{应用示例}：已知 SO\textsubscript{2} 氧化和 NO\textsubscript{2} 分解的 \(K^{\circ}\)，可求 SO\textsubscript{2} + NO\textsubscript{2} → SO\textsubscript{3} + NO 的 \(K^{\circ}\)。这在处理耦合平衡(如沉淀-配位、酸碱-配位)时是关键工具。

\begin{center}\rule{0.5\linewidth}{0.5pt}\end{center}

\subsection{四、反应方向判断：Q 与 K 的比较}\label{ux56dbux53cdux5e94ux65b9ux5411ux5224ux65adq-ux4e0e-k-ux7684ux6bd4ux8f83}

\subsubsection{4.1 Q 的书写}\label{q-ux7684ux4e66ux5199}

Q 的表达式与 \(K^{\circ}\) 形式完全相同，区别在于 Q 用\textbf{任意时刻}的数据，K 用\textbf{平衡时}的数据。

{\def\LTcaptype{none} % do not increment counter
\begin{longtable}[]{@{}
  >{\raggedright\arraybackslash}p{(\linewidth - 4\tabcolsep) * \real{0.3333}}
  >{\raggedright\arraybackslash}p{(\linewidth - 4\tabcolsep) * \real{0.3333}}
  >{\raggedright\arraybackslash}p{(\linewidth - 4\tabcolsep) * \real{0.3333}}@{}}
\toprule\noalign{}
\begin{minipage}[b]{\linewidth}\raggedright
反应类型
\end{minipage} & \begin{minipage}[b]{\linewidth}\raggedright
示例
\end{minipage} & \begin{minipage}[b]{\linewidth}\raggedright
Q 表达式
\end{minipage} \\
\midrule\noalign{}
\endhead
\bottomrule\noalign{}
\endlastfoot
气相 & \(\mathrm{N_2 + 3H_2 \rightleftharpoons 2NH_3}\) & \(Q = \frac{(p_{\mathrm{NH_3}}/p^\theta)^2}{(p_{\mathrm{N_2}}/p^\theta)(p_{\mathrm{H_2}}/p^\theta)^3}\) \\
液相 & \(\mathrm{HAc \rightleftharpoons H^+ + Ac^-}\) & \(Q = \frac{[\mathrm{H^+}][\mathrm{Ac^-}]}{[\mathrm{HAc}]}\) \\
多相 & \(\mathrm{CaCO_3(s) \rightleftharpoons CaO(s) + CO_2(g)}\) & \(Q = p_{\mathrm{CO_2}}/p^\theta\) \\
\end{longtable}
}

\subsubsection{4.2 Q-K 判断的定量分析}\label{q-k-ux5224ux65adux7684ux5b9aux91cfux5206ux6790}

以 H\textsubscript{2} + I\textsubscript{2} ⇌ 2HI 体系为例(713 K，K = 50.3)：

{\def\LTcaptype{none} % do not increment counter
\begin{longtable}[]{@{}ccccccl@{}}
\toprule\noalign{}
状态 & (H\textsubscript{2}) & (I\textsubscript{2}) & (HI) & Q & Q vs K & 方向 \\
\midrule\noalign{}
\endhead
\bottomrule\noalign{}
\endlastfoot
I & 1.00 & 1.00 & 1.00 & 1.00 & Q \textless{} K & 正向自发 \\
II & 1.00 & 1.00 & 0.001 & 10\textsuperscript{-6} & Q ≪ K & 正向自发(程度大) \\
III & 0.22 & 0.22 & 1.56 & 50 & Q ≈ K & 平衡 \\
IV & 0.22 & 0.22 & 2.56 & 140 & Q \textgreater{} K & 逆向自发 \\
V & 1.22 & 0.22 & 1.56 & 9.1 & Q \textless{} K & 正向自发 \\
\end{longtable}
}

\textbf{核心思想}：Q/K 之值不仅决定反应进行的方向，而且也表明了起始状态和平衡状态之间的差距，也就预示了平衡移动的相对多少。

\begin{center}\rule{0.5\linewidth}{0.5pt}\end{center}

\subsection{五、平衡移动}\label{ux4e94ux5e73ux8861ux79fbux52a8}

\subsubsection{5.1 勒夏特列原理(定性)}\label{ux52d2ux590fux7279ux5217ux539fux7406ux5b9aux6027}

改变平衡体系的条件之一，平衡就向着\textbf{减弱这个改变}的方向移动。

{\def\LTcaptype{none} % do not increment counter
\begin{longtable}[]{@{}lll@{}}
\toprule\noalign{}
条件变化 & 平衡移动方向 & Q-K 分析 \\
\midrule\noalign{}
\endhead
\bottomrule\noalign{}
\endlastfoot
增加反应物浓度 & 正向移动 & 分母增大 → Q \textless{} K \\
减少生成物浓度 & 正向移动 & 分子减小 → Q \textless{} K \\
增大总压(Δν \textless{} 0) & 向分子数减少方向 & Q \textless{} K \\
减小总压(Δν \textgreater{} 0) & 向分子数增加方向 & Q \textgreater{} K \\
升温(正反应吸热) & 正向移动 & K 增大 \\
升温(正反应放热) & 逆向移动 & K 减小 \\
加催化剂 & \textbf{不移动} & 正逆速率同等加快，K 不变 \\
\end{longtable}
}

\textbf{口诀的边界}：勒夏特列原理是定性近似。在复杂条件(如恒压充惰性气体、多因素同时变化)下，必须用 Q-K 定量分析，不能死记口诀。

\subsubsection{5.2 浓度与压力的影响---改变 Q}\label{ux6d53ux5ea6ux4e0eux538bux529bux7684ux5f71ux54cdux6539ux53d8-q}

\textbf{浓度变化}：增加反应物浓度 → 分母增大 → Q \textless{} K → 正向移动。

\textbf{恒温改变总压}：
- Δν \textless{} 0：增压 → 正向移动(分子数减少方向)
- Δν \textgreater{} 0：增压 → 逆向移动
- Δν = 0：压力变化\textbf{不影响}平衡

\textbf{实例：N\textsubscript{2}O\textsubscript{4} 分解的压力效应}

在 325 K、100 kPa 下，N\textsubscript{2}O\textsubscript{4}(g) ⇌ 2NO\textsubscript{2}(g) 的分解率 α = 50.2\%。若压力增至 1000 kPa：

设起始 1 mol N\textsubscript{2}O\textsubscript{4}，分解率 α'：

\[
$K_{\1}$ = \frac{(2\alpha'/n_{总})^2 \cdot p^2}{((1-\alpha')/n_{总}) \cdot p} = \frac{4\alpha'^2 p}{(1-\alpha'^2)(1+\alpha')}
\]

代入 \(K_{\1}\) 值(由 100 kPa 数据求得)和 p = 1000 kPa，解得 α' ≪ 50.2\%。\textbf{增压使分解率降低}，符合''增压向分子数减少方向''的预测。

\subsubsection{5.3 惰性气体的影响(竞赛高频)}\label{ux60f0ux6027ux6c14ux4f53ux7684ux5f71ux54cdux7adeux8d5bux9ad8ux9891}

{\def\LTcaptype{none} % do not increment counter
\begin{longtable}[]{@{}
  >{\raggedright\arraybackslash}p{(\linewidth - 4\tabcolsep) * \real{0.3333}}
  >{\raggedright\arraybackslash}p{(\linewidth - 4\tabcolsep) * \real{0.3333}}
  >{\raggedright\arraybackslash}p{(\linewidth - 4\tabcolsep) * \real{0.3333}}@{}}
\toprule\noalign{}
\begin{minipage}[b]{\linewidth}\raggedright
条件
\end{minipage} & \begin{minipage}[b]{\linewidth}\raggedright
对平衡的影响
\end{minipage} & \begin{minipage}[b]{\linewidth}\raggedright
原因
\end{minipage} \\
\midrule\noalign{}
\endhead
\bottomrule\noalign{}
\endlastfoot
\textbf{恒容}充惰气 & \textbf{不移动} & 各组分分压不变(\(p_i = n_iRT/V\))，Q 不变 \\
\textbf{恒压}充惰气 & \textbf{向分子数增加方向移动} & 体积增大，各组分分压降低，等效减压 \\
\end{longtable}
}

\textbf{最常见错误}：认为''恒容充惰气总压增大，平衡就移动''。实际上各组分分压没变，Q 没变，平衡不移动！

\textbf{PCl\textsubscript{5} 解离度分析}(经典辨析题)：

PCl\textsubscript{5}(g) ⇌ PCl\textsubscript{3}(g) + Cl\textsubscript{2}(g)，分析以下条件下解离度 α 的变化：

{\def\LTcaptype{none} % do not increment counter
\begin{longtable}[]{@{}lcl@{}}
\toprule\noalign{}
条件 & α 的变化 & 分析 \\
\midrule\noalign{}
\endhead
\bottomrule\noalign{}
\endlastfoot
降压(扩大体积) & 增大 & 等效减压，向分子数增加方向 \\
恒压充 N\textsubscript{2} & 增大 & 体积增大，等效减压 \\
恒容充 N\textsubscript{2} & 不变 & 各组分分压不变，Q 不变 \\
恒压充 Cl\textsubscript{2} & 减小 & 增加产物浓度 \\
恒容充 Cl\textsubscript{2} & 减小 & 增加产物浓度 \\
\end{longtable}
}

\subsubsection{5.4 温度的影响---改变 K}\label{ux6e29ux5ea6ux7684ux5f71ux54cdux6539ux53d8-k}

温度是唯一能改变 K 的因素。\textbf{van't Hoff 方程}：

\[
\boxed{\ln\frac{K_2}{K_1} = \frac{\Delta_rH^\circ}{R}\left(\frac{1}{T_1} - \frac{1}{T_2}\right)}
\]

{\def\LTcaptype{none} % do not increment counter
\begin{longtable}[]{@{}lll@{}}
\toprule\noalign{}
反应热效应 & 升温效果 & 降温效果 \\
\midrule\noalign{}
\endhead
\bottomrule\noalign{}
\endlastfoot
吸热(Δ\_rH° \textgreater{} 0) & K 增大，正向移动 & K 减小，逆向移动 \\
放热(Δ\_rH° \textless{} 0) & K 减小，逆向移动 & K 增大，正向移动 \\
\end{longtable}
}

\textbf{记忆}：升温有利于\textbf{吸热方向}。

\textbf{van't Hoff 方程的线性形式}：

\[
\ln K = -\frac{\Delta_rH^\circ}{RT} + \frac{\Delta_rS^\circ}{R}
\]

以 ln K 对 1/T 作图，斜率 = −Δ\_rH°/R，截距 = Δ\_rS°/R。这是实验测定 Δ\_rH° 的重要方法。

\subsubsection{5.5 Clausius-Clapeyron 方程(相变平衡特例)}\label{clausius-clapeyron-ux65b9ux7a0bux76f8ux53d8ux5e73ux8861ux7279ux4f8b}

液-气或固-气相变可视为平衡的特例。对相变过程：

\[
\mathrm{Br_2(l) \rightleftharpoons Br_2(g)}
\]

平衡时蒸气压等于外压，ΔG = 0。将 van't Hoff 方程应用于相变，得到 \textbf{Clausius-Clapeyron 方程}：

\[
\boxed{\ln\frac{p_2}{p_1} = -\frac{\Delta_rH_{vap}}{R}\left(\frac{1}{T_2} - \frac{1}{T_1}\right)}
\]

其中 Δ\_rH\_vap 为摩尔蒸发热(近似视为常数)。

\textbf{应用1：由两组蒸气压数据求蒸发热}

异丙醇在 2.4°C(275.6 K)时蒸气压为 1.33 kPa，39.5°C(312.7 K)时为 13.3 kPa：

\[
\ln\frac{13.3}{1.33} = -\frac{\Delta_rH_{vap}}{8.314}\left(\frac{1}{312.7} - \frac{1}{275.6}\right)
\]

\[
2.303 = -\frac{\Delta_rH_{vap}}{8.314} \times (-4.38 \times 10^{-4})
\]

\[
\Delta_rH_{vap} = 43.5\ \mathrm{kJ/mol}
\]

\textbf{应用2：估算正常沸点}

正常沸点定义：蒸气压 = 100 kPa(1 bar)时的温度。代入上式求 T\textsubscript{2}：

\[
\ln\frac{100}{1.33} = \frac{43500}{8.314}\left(\frac{1}{275.6} - \frac{1}{T_b}\right)
\]

解得 \(T_b \approx 362\ \mathrm{K}\)(89°C)，与文献值 82.6°C 接近(近似处理的误差)。

\textbf{应用3：估算任意温度下的蒸气压}：已知蒸发热和一个温度点的蒸气压，可估算其他温度下的蒸气压。这在相图分析和蒸馏计算中非常有用。

\subsubsection{5.6 恒压 vs 恒容升温对平衡移动程度的比较}\label{ux6052ux538b-vs-ux6052ux5bb9ux5347ux6e29ux5bf9ux5e73ux8861ux79fbux52a8ux7a0bux5ea6ux7684ux6bd4ux8f83}

对于 N\textsubscript{2}O\textsubscript{4} ⇌ 2NO\textsubscript{2}(吸热、气体分子数增加)：

{\def\LTcaptype{none} % do not increment counter
\begin{longtable}[]{@{}
  >{\raggedright\arraybackslash}p{(\linewidth - 6\tabcolsep) * \real{0.2500}}
  >{\raggedright\arraybackslash}p{(\linewidth - 6\tabcolsep) * \real{0.2500}}
  >{\raggedright\arraybackslash}p{(\linewidth - 6\tabcolsep) * \real{0.2500}}
  >{\raggedright\arraybackslash}p{(\linewidth - 6\tabcolsep) * \real{0.2500}}@{}}
\toprule\noalign{}
\begin{minipage}[b]{\linewidth}\raggedright
条件
\end{minipage} & \begin{minipage}[b]{\linewidth}\raggedright
升温效应
\end{minipage} & \begin{minipage}[b]{\linewidth}\raggedright
附加效应
\end{minipage} & \begin{minipage}[b]{\linewidth}\raggedright
净效果
\end{minipage} \\
\midrule\noalign{}
\endhead
\bottomrule\noalign{}
\endlastfoot
\textbf{恒压}升温 & K 增大，右移 & 无附加效应 & 移动程度大 \\
\textbf{恒容}升温 & K 增大，右移 & 温度升高→压强增大→左移(部分抵消) & 移动程度小 \\
\end{longtable}
}

\textbf{结论}：恒压条件下升温对平衡移动的程度大于恒容条件。这是因为恒容升温时，温度效应(右移)和压强效应(左移)部分抵消。这是国初真题中常考的辨析点。

\begin{center}\rule{0.5\linewidth}{0.5pt}\end{center}

\subsection{六、平衡计算}\label{ux516dux5e73ux8861ux8ba1ux7b97}

\subsubsection{6.1 ICE 表方法}\label{ice-ux8868ux65b9ux6cd5}

\textbf{ICE = Initial(初始)→ Change(变化)→ Equilibrium(平衡)}

以 \(a\mathrm{A} + b\mathrm{B} \rightleftharpoons d\mathrm{D} + e\mathrm{E}\) 为例：

{\def\LTcaptype{none} % do not increment counter
\begin{longtable}[]{@{}lcccc@{}}
\toprule\noalign{}
& A & B & D & E \\
\midrule\noalign{}
\endhead
\bottomrule\noalign{}
\endlastfoot
I(初始) & \(c_{A,0}\) & \(c_{B,0}\) & \(c_{D,0}\) & \(c_{E,0}\) \\
C(变化) & \(-ax\) & \(-bx\) & \(+dx\) & \(+ex\) \\
E(平衡) & \(c_{A,0}-ax\) & \(c_{B,0}-bx\) & \(c_{D,0}+dx\) & \(c_{E,0}+ex\) \\
\end{longtable}
}

代入 K 表达式求解 \(x\)。

\subsubsection{6.2 平衡转化率}\label{ux5e73ux8861ux8f6cux5316ux7387}

\[
\alpha = \frac{\text{已转化的量}}{\text{起始总量}} \times 100\%
\]

\textbf{关键区分}：
- K 与起始状态\textbf{无关}(只与温度有关)
- α 与起始状态\textbf{密切相关}(不同起始条件，同一温度下 α 不同)
- 多个反应物时，各反应物的 α 通常不同

\subsubsection{6.3 分解压}\label{ux5206ux89e3ux538b}

纯固体分解产生气体时，平衡气体的总压力称为分解压。

\begin{itemize}
\tightlist
\item
  单一气体产物(如 CaCO\textsubscript{3})：分解压 = \(p_{\mathrm{CO_2}}\) = \(K_{\1}\)
\item
  多种气体产物(如 NH\textsubscript{4}HS)：分解压 = \(p_{\mathrm{NH_3}} + p_{\mathrm{H_2S}}\)
\end{itemize}

\textbf{实例：NH\textsubscript{4}HS 分解平衡}

NH\textsubscript{4}HS(s) ⇌ NH\textsubscript{3}(g) + H\textsubscript{2}S(g)，500 K 时 \(K_{\1}\) = 0.108 (atm)\textsuperscript{2}。

\begin{enumerate}
\def\labelenumi{(\alph{enumi})}
\tightlist
\item
  通入 0.50 atm NH\textsubscript{3} 和 0.50 atm H\textsubscript{2}S，判断是否平衡：
\end{enumerate}

\[
Q = p(\mathrm{NH_3}) \times p(\mathrm{H_2S}) = 0.50 \times 0.50 = 0.25 > $K_{\1}$ = 0.108
\]

Q \textgreater{} K，不平衡。逆向自发，NH\textsubscript{4}HS(s) 生成(气体结合沉积为固体)。

\begin{enumerate}
\def\labelenumi{(\alph{enumi})}
\setcounter{enumi}{1}
\tightlist
\item
  求新平衡时各组分分压：设平衡时 H\textsubscript{2}S 分压为 x atm，则 NH\textsubscript{3} 分压也为 x atm(1:1 产气)，总压 = 2x。但注意通入的 NH\textsubscript{3} 过量，设平衡时 p(H\textsubscript{2}S) = y：
\end{enumerate}

\[
$K_{\1}$ = p(\mathrm{NH_3}) \times p(\mathrm{H_2S}) = (0.50 - \Delta)(0.50 - \Delta)
\]

其中 Δ 为逆向反应消耗的量。由 \(K_{\1}\) = (0.50 − Δ)\textsuperscript{2} = 0.108，解得 Δ = 0.16 atm，平衡时 p(H\textsubscript{2}S) = p(NH\textsubscript{3}) = 0.34 atm。

\subsubsection{6.4 平衡计算通用流程}\label{ux5e73ux8861ux8ba1ux7b97ux901aux7528ux6d41ux7a0b}

\begin{verbatim}
1. 写出配平方程式
2. 写出 K 表达式
3. 列 ICE 表
4. 代入 K 表达式
5. 求解(近似法 or 精确法)
6. 验证近似合理性
\end{verbatim}

\textbf{近似法条件}：当 K 很小(如 \textless{} 10\textsuperscript{-4})时，可假设反应物转化量远小于起始量，简化计算。但需最后验证假设是否合理。

\begin{center}\rule{0.5\linewidth}{0.5pt}\end{center}

\subsection{七、知识网络与方法论}\label{ux4e03ux77e5ux8bc6ux7f51ux7edcux4e0eux65b9ux6cd5ux8bba}

\subsubsection{7.1 ``Q-K'' 分析框架}\label{q-k-ux5206ux6790ux6846ux67b6}

\begin{verbatim}
改变浓度或压力？
  → Q 改变，K 不变
  → 比较 Q 与 K 判断方向

改变温度？
  → K 改变，Q 不变
  → 用 van't Hoff 方程计算新 K
  → 再比较 Q 与 K
\end{verbatim}

\subsubsection{7.2 三大核心公式的关系}\label{ux4e09ux5927ux6838ux5fc3ux516cux5f0fux7684ux5173ux7cfb}

\begin{verbatim}
$\Delta \1^{\circ}$ = -RT ln $K^{\circ}$     (热力学→平衡常数)
ΔG = $\Delta \1^{\circ}$ + RT ln Q  (任意状态→方向判据)
ΔG = RT ln(Q/$K^{\circ}$)    (综合判据)
\end{verbatim}

\subsubsection{7.3 温度影响的两种处理}\label{ux6e29ux5ea6ux5f71ux54cdux7684ux4e24ux79cdux5904ux7406}

{\def\LTcaptype{none} % do not increment counter
\begin{longtable}[]{@{}
  >{\raggedright\arraybackslash}p{(\linewidth - 4\tabcolsep) * \real{0.3333}}
  >{\raggedright\arraybackslash}p{(\linewidth - 4\tabcolsep) * \real{0.3333}}
  >{\raggedright\arraybackslash}p{(\linewidth - 4\tabcolsep) * \real{0.3333}}@{}}
\toprule\noalign{}
\begin{minipage}[b]{\linewidth}\raggedright
场景
\end{minipage} & \begin{minipage}[b]{\linewidth}\raggedright
工具
\end{minipage} & \begin{minipage}[b]{\linewidth}\raggedright
公式
\end{minipage} \\
\midrule\noalign{}
\endhead
\bottomrule\noalign{}
\endlastfoot
已知两个温度的 K & van't Hoff 两点法 & \(\ln(K_2/K_1) = (\Delta_rH°/R)(1/T_1 - 1/T_2)\) \\
已知蒸气压数据 & C-C 方程 & \(\ln(p_2/p_1) = (-\Delta_rH_{vap}/R)(1/T_2 - 1/T_1)\) \\
\end{longtable}
}

\begin{center}\rule{0.5\linewidth}{0.5pt}\end{center}

\subsection{八、典型例题}\label{ux516bux5178ux578bux4f8bux9898}

\textbf{题干}：将 0.100 mol N\textsubscript{2}O\textsubscript{4} 充入 1 dm\textsuperscript{3} 烧瓶，373 K 下平衡后测得 {[}N\textsubscript{2}O\textsubscript{4}{]} = 0.040 mol/L。求 \(K_{\1}\)。

\textbf{解析}：

{\def\LTcaptype{none} % do not increment counter
\begin{longtable}[]{@{}lcc@{}}
\toprule\noalign{}
& N\textsubscript{2}O\textsubscript{4} & 2NO\textsubscript{2} \\
\midrule\noalign{}
\endhead
\bottomrule\noalign{}
\endlastfoot
起始 & 0.100 & 0 \\
变化 & −0.060 & +0.120 \\
平衡 & 0.040 & 0.120 \\
\end{longtable}
}

\[K_c = \frac{[\mathrm{NO_2}]^2}{[\mathrm{N_2O_4}]} = \frac{(0.120)^2}{0.040} = 0.36\]

\textbf{答案}：\(K_{\1}\) = 0.36(373 K)。

从不同起始条件(如只充 NO\textsubscript{2})出发，最终平衡时 {[}NO\textsubscript{2}{]}\textsuperscript{2}/{[}N\textsubscript{2}O\textsubscript{4}{]} 的比值相同，这正是''平衡常数只与温度有关''的实验验证。

\begin{center}\rule{0.5\linewidth}{0.5pt}\end{center}

\textbf{题干}：某反应 A(s) ⇌ B(s) + C(g)，\(\Delta \1^{\circ}\)(298 K) = +40.0 kJ/mol。(1) 求 \(K_{\1}\)°；(2) 当 p\_C = 1.0 Pa 时，反应能否正向自发？

\textbf{解析}：

\begin{enumerate}
\def\labelenumi{(\arabic{enumi})}
\tightlist
\item
  由 \(\Delta \1^{\circ}\) = −2.30RT lg \(K_{\1}\)°：
\end{enumerate}

\[\lg $K_{\1}$ = -\frac{40000}{2.30 \times 8.314 \times 298} = -7.02\]

\[K_p^\theta = 9.5 \times 10^{-8}\]

\begin{enumerate}
\def\labelenumi{(\arabic{enumi})}
\setcounter{enumi}{1}
\tightlist
\item
  Q\_p = p\_C/p° = 1.0 × 10\textsuperscript{-5} / 1.0 × 10\textsuperscript{2} = 1.0 × 10\textsuperscript{-7}
\end{enumerate}

\[\Delta G = 2.30RT\lg\frac{Q_p}{$K_{\1}$} = 2.30 \times 8.314 \times 10^{-3} \times 298 \times \lg\frac{1.0 \times 10^{-7}}{9.5 \times 10^{-8}} = +5.7\ \mathrm{kJ/mol} > 0\]

\textbf{答案}：(1) \(K_{\1}\)° = 9.5 × 10\textsuperscript{-8}；(2) ΔG \textgreater{} 0，反应不能正向自发。即使 p\_C 从标态(100 kPa)降低到 1.0 Pa(降低 5 个量级)，ΔG 仍为正值。

\textbf{启示}：当 \(\Delta \1^{\circ}\) 的绝对值相当大时(如 \textgreater{} 40 kJ/mol)，Q 的变化很难改变 ΔG 的正负号。此时用 \(\Delta \1^{\circ}\) 粗估反应方向是可行的。

\begin{center}\rule{0.5\linewidth}{0.5pt}\end{center}

\textbf{题干}：合成氨 N\textsubscript{2}(g) + 3H\textsubscript{2}(g) ⇌ 2NH\textsubscript{3}(g) 在 673 K 时 K\_p1° = 1.6 × 10\textsuperscript{-4}，Δ\_rH° = −92.4 kJ/mol(视为常数)。求 773 K 时的 K\_p2°。

\textbf{解析}：

\[\ln\frac{K_2}{K_1} = \frac{\Delta_rH^\circ}{R}\left(\frac{1}{T_1} - \frac{1}{T_2}\right) = \frac{-92400}{8.314}\left(\frac{1}{673} - \frac{1}{773}\right)\]

\[= -11114 \times 1.93 \times 10^{-4} = -2.14\]

\[K_2 = K_1 \times e^{-2.14} = 1.6 \times 10^{-4} \times 0.118 = 1.9 \times 10^{-5}\]

\textbf{答案}：K\_p2° = 1.9 × 10\textsuperscript{-5}(升温使 K 减小，符合放热反应规律)。

\begin{center}\rule{0.5\linewidth}{0.5pt}\end{center}

\textbf{题干}：1000°C、3000 kPa 下 CO\textsubscript{2}(g) + C(s) ⇌ 2CO(g) 平衡时 CO\textsubscript{2} 摩尔分数为 0.17。求：
(1) \(K^{\circ}\)
(2) 若总压降至 2000 kPa，新平衡 CO\textsubscript{2} 摩尔分数
(3) 恒容充 N\textsubscript{2} 至 4000 kPa，平衡是否移动

\textbf{解析}：

\begin{enumerate}
\def\labelenumi{(\arabic{enumi})}
\tightlist
\item
  p(CO) = 3000 × 0.83 = 2490 kPa，p(CO\textsubscript{2}) = 3000 × 0.17 = 510 kPa
\end{enumerate}

\[K^\theta = \frac{(2490/100)^2}{510/100} = 122\]

\begin{enumerate}
\def\labelenumi{(\arabic{enumi})}
\setcounter{enumi}{1}
\tightlist
\item
  设新平衡 CO\textsubscript{2} 摩尔分数为 x：
\end{enumerate}

\[K^\theta = \frac{[2000(1-x)/100]^2}{2000x/100} = 122\]

\[\frac{400(1-x)^2}{20x} = 122 \implies 20x^2 - 162x + 20 = 0\]

\[x = \frac{162 - \sqrt{162^2 - 4 \times 20 \times 20}}{40} = 0.129\]

\begin{enumerate}
\def\labelenumi{(\arabic{enumi})}
\setcounter{enumi}{2}
\tightlist
\item
  恒容充惰气：各组分分压不变 → Q 不变 → \textbf{平衡不移动}。
\end{enumerate}

\textbf{答案}：(1) \(K^{\circ}\) = 122；(2) CO\textsubscript{2} 摩尔分数 = 0.129；(3) 不移动。

\begin{center}\rule{0.5\linewidth}{0.5pt}\end{center}

\textbf{题干}：已知 298 K 时：
- AgCl(s) ⇌ Ag\textsuperscript{+}(aq) + Cl\textsuperscript{-}(aq)，\(K_{\1}\)° = 1.8 × 10\textsuperscript{-10}
- Ag\textsuperscript{+}(aq) + 2NH\textsubscript{3}(aq) ⇌ Ag(NH\textsubscript{3})\textsubscript{2}\textsuperscript{+}(aq)，\(K_{\1}\)° = 1.7 × 10\textsuperscript{7}

求 AgCl(s) + 2NH\textsubscript{3}(aq) ⇌ Ag(NH\textsubscript{3})\textsubscript{2}\textsuperscript{+}(aq) + Cl\textsuperscript{-}(aq) 的 \(K^{\circ}\)。解释为什么 AgCl 能溶于浓氨水。

\textbf{解析}：

目标反应 = AgCl 溶解 + 配位，由多重平衡规则：

\[K^\theta = K_{sp}^\theta \times $K_{\1}$^\theta = 1.8 \times 10^{-10} \times 1.7 \times 10^7 = 3.1 \times 10^{-3}\]

\(K^{\circ}\) 较小但不极小。在浓氨水({[}NH\textsubscript{3}{]} ≈ 15 mol/L)中：

\[Q = \frac{[\mathrm{Ag(NH_3)_2^+}][\mathrm{Cl^-}]}{[\mathrm{NH_3}]^2} \approx \frac{s \cdot s}{15^2} = \frac{s^2}{225}\]

Q ≪ K，反应正向进行，AgCl 能溶解。在稀氨水中 {[}NH\textsubscript{3}{]} 不够大，Q \textgreater{} K，溶解不完全。

\textbf{答案}：\(K^{\circ}\) = 3.1 × 10\textsuperscript{-3}。浓氨水中 AgCl 能溶解，稀氨水中溶解不完全。

这是定性分析中''银的氯化物溶于浓氨水''的经典解释，也是沉淀-配位耦合平衡的典型应用。

\begin{center}\rule{0.5\linewidth}{0.5pt}\end{center}

\subsection{九、本节总结速查}\label{ux4e5dux672cux8282ux603bux7ed3ux901fux67e5}

{\def\LTcaptype{none} % do not increment counter
\begin{longtable}[]{@{}
  >{\raggedright\arraybackslash}p{(\linewidth - 4\tabcolsep) * \real{0.3333}}
  >{\raggedright\arraybackslash}p{(\linewidth - 4\tabcolsep) * \real{0.3333}}
  >{\raggedright\arraybackslash}p{(\linewidth - 4\tabcolsep) * \real{0.3333}}@{}}
\toprule\noalign{}
\begin{minipage}[b]{\linewidth}\raggedright
核心概念
\end{minipage} & \begin{minipage}[b]{\linewidth}\raggedright
关键公式
\end{minipage} & \begin{minipage}[b]{\linewidth}\raggedright
注意点
\end{minipage} \\
\midrule\noalign{}
\endhead
\bottomrule\noalign{}
\endlastfoot
\(K_{\1}\) & \(K_c = \frac{[\mathrm{D}]^d[\mathrm{E}]^e}{[\mathrm{A}]^a[\mathrm{B}]^b}\) & 纯固体/液体不写入 \\
\(K_{\1}\) = \(K_{\1}\)(RT)\^{}Δν & Δν = 产物系数和 − 反应物系数和 & R 值需与单位匹配 \\
标准平衡常数 \(K^{\circ}\) & 无量纲，\(\Delta_rG^\circ = -RT\ln K^\theta\) & c°=1 mol/L，p°=100 kPa \\
多重平衡 & 反应相加 → K 相乘 & 竞赛中处理耦合平衡的关键 \\
Q vs K & Q \textless{} K 正向，Q \textgreater{} K 逆向 & Q 用任意态，K 用平衡态 \\
浓度/压力影响 & 改变 Q，K 不变 & 恒容充惰气不影响平衡 \\
温度影响 & 改变 K，Q 不变 & van't Hoff 方程 \\
C-C 方程 & 相变平衡的 van't Hoff 特例 & 蒸发热/沸点计算 \\
ICE 表 & I→C→E 三行 & Δν=0 可开方，否则解二次 \\
\end{longtable}
}

\begin{center}\rule{0.5\linewidth}{0.5pt}\end{center}

\subsection{十、三级练习题}\label{ux5341ux4e09ux7ea7ux7ec3ux4e60ux9898}

\textbf{1.} 写出下列反应的 \(K_{\1}\) 和 \(K_{\1}\) 表达式(如适用)：
(a) \(\mathrm{N_2(g) + 3H_2(g) \rightleftharpoons 2NH_3(g)}\)
(b) \(\mathrm{CaCO_3(s) \rightleftharpoons CaO(s) + CO_2(g)}\)
(c) \(\mathrm{CH_3COOH(aq) \rightleftharpoons H^+(aq) + CH_3COO^-(aq)}\)

\textbf{2.} 已知 298 K 时某反应的 Δ\_rG° = +20.0 kJ/mol，计算 \(K^{\circ}\)。反应在标准状态下正向能否自发进行？

\textbf{3.} 判断下列说法的正误：
(a) 平衡常数越大，反应速率越快。
(b) 加入催化剂会改变平衡常数。
(c) 恒容条件下充入惰性气体，平衡不移动。
(d) 放热反应升温后平衡常数增大。

\textbf{4.} 对反应 \(\mathrm{2SO_2(g) + O_2(g) \rightleftharpoons 2SO_3(g)}\)，Δν = −1。增大总压平衡向哪个方向移动？为什么？

\textbf{5.} 某温度下 \(K_c = 4\)，反应 \(\mathrm{A(g) \rightleftharpoons 2B(g)}\) 的 A 起始浓度为 1.0 mol/L，求平衡时 A 和 B 的浓度。

\textbf{6.} 简述''恒容充惰气''与''恒压充惰气''对化学平衡影响的区别。

\textbf{7.} 已知反应(1)的 K\textsubscript{1} 和反应(2)的 K\textsubscript{2}，若反应(3) = 2×(1) − (2)，求 K\textsubscript{3} 的表达式。

\textbf{8.} 某放热反应在 500 K 时 \(K^{\circ}\) = 100。升温到 600 K 后 \(K^{\circ}\) 是增大还是减小？用 van't Hoff 原理说明。

\textbf{9.} 在恒温恒容条件下，向已达平衡的 \(\mathrm{N_2O_4 \rightleftharpoons 2NO_2}\) 体系中加入 NO\textsubscript{2}，平衡向哪个方向移动？从 Q-K 角度解释。

\textbf{10.} 某分解反应 \(\mathrm{A(s) \rightleftharpoons B(g) + C(g)}\) 在 500 K 时 \(K_{\1}\) = 0.04 (kPa)\textsuperscript{2}。求该温度下 A 的分解压。

\begin{center}\rule{0.5\linewidth}{0.5pt}\end{center}

\textbf{11.} \textbf{(ICE 表综合)} 某温度下 \(\mathrm{H_2(g) + I_2(g) \rightleftharpoons 2HI(g)}\) 的 \(K_c = 49\)。在 10 L 容器中充入 1 mol H\textsubscript{2} 和 1 mol I\textsubscript{2}，求平衡时各物质的浓度和 H\textsubscript{2} 的转化率。

\textbf{12.} \textbf{(\(K_{\1}\) 与 \(K_{\1}\) 转换)} 727°C 时 \(\mathrm{FeO(s) + CO(g) \rightleftharpoons Fe(s) + CO_2(g)}\) 的 \(K_p = 0.681\)。求该温度下的 \(K_{\1}\)，并说明对于 Δν = 0 的反应 \(K_{\1}\) 与 \(K_{\1}\) 的关系。

\textbf{13.} \textbf{(压力对平衡的影响)} 在 2000 K、100 kPa 下，反应 \(\mathrm{N_2(g) + O_2(g) \rightleftharpoons 2NO(g)}\) 的 \(K^{\circ}\) = 4.08 × 10\textsuperscript{-4}。若总压增至 1000 kPa(温度不变)，NO 的摩尔分数如何变化？通过计算说明。

\textbf{14.} \textbf{(多重平衡)} 已知 298 K 时：
- \(\mathrm{Ag^+(aq) + 2NH_3(aq) \rightleftharpoons Ag(NH_3)_2^+(aq)}\)，\(K_f^\theta = 1.7 \times 10^7\)
- \(\mathrm{AgCl(s) \rightleftharpoons Ag^+(aq) + Cl^-(aq)}\)，\(K_{sp}^\theta = 1.8 \times 10^{-10}\)

求 \(\mathrm{AgCl(s) + 2NH_3(aq) \rightleftharpoons Ag(NH_3)_2^+(aq) + Cl^-(aq)}\) 的 \(K^{\circ}\)，并解释为什么 AgCl 能溶于浓氨水。

\textbf{15.} \textbf{(van't Hoff 定量)} 已知 298 K 时反应 \(\mathrm{N_2O_4(g) \rightleftharpoons 2NO_2(g)}\) 的 Δ\_rH° = +57.2 kJ/mol，\(K^{\circ}\) = 0.148。估算 \(K^{\circ}\) = 1 时的温度。

\textbf{16.} \textbf{(分解压应用)} 500 K 时 \(\mathrm{NH_4HS(s) \rightleftharpoons NH_3(g) + H_2S(g)}\) 的 \(K_{\1}\) = 0.108 (atm)\textsuperscript{2}。在密闭容器中通入 0.50 atm NH\textsubscript{3} 和 H\textsubscript{2}S 各 0.50 atm，判断：
(a) 反应是否已达平衡？
(b) 若不平衡，NH\textsubscript{4}HS(s) 会分解还是生成？

\textbf{17.} \textbf{(转化率与起始条件)} 对反应 \(\mathrm{A(g) + B(g) \rightleftharpoons 2C(g)}\)，\(K_{\1}\) = 4.0。比较以下两种起始条件下 A 的转化率：
(a) A 和 B 各 1.0 mol/L
(b) A 为 2.0 mol/L，B 为 0.5 mol/L

\textbf{18.} \textbf{(工业合成氨条件分析)} 合成氨 N\textsubscript{2} + 3H\textsubscript{2} ⇌ 2NH\textsubscript{3}，Δ\_rH° = −92.4 kJ/mol。
(a) 从热力学角度分析：低温有利于 K 增大，为什么工业上不用极低温度？
(b) 从动力学角度解释：为什么需要用铁催化剂和 400−500°C？
(c) 从平衡角度解释：为什么工业上采用 200−300 atm 高压？

\textbf{19.} \textbf{(耦合平衡)} 在 0.10 mol/L FeCl\textsubscript{3} 溶液中通入 H\textsubscript{2}S 至饱和({[}H\textsubscript{2}S{]} = 0.10 mol/L)，问：
(a) \(\mathrm{Fe^{3+}}\) 能否将 \(\mathrm{H_2S}\) 氧化为 \(\mathrm{S}\)？(已知 \(\mathrm{Fe^{3+}/Fe^{2+}}\) 的 \(E^{\circ}\) = 0.77 V，\(\mathrm{S/H_2S}\) 的 \(E^{\circ}\) = 0.14 V)
(b) 若能反应，计算该反应的 \(K^{\circ}\)。
(c) 反应达平衡后 \(\mathrm{Fe^{3+}}\) 的浓度为多少？(已知 \(K_{sp}^\theta(\mathrm{Fe_2S_3}) = 1.0 \times 10^{-88}\)，讨论 Fe\textsubscript{2}S\textsubscript{3} 能否沉淀)

\begin{center}\rule{0.5\linewidth}{0.5pt}\end{center}

\textbf{20.} \textbf{(全国初赛)} 某温度下，反应 \(\mathrm{CO(g) + H_2O(g) \rightleftharpoons CO_2(g) + H_2(g)}\) 在密闭容器中达到平衡。已知 \(K_{\1}\) = 5.0，起始时 c(CO) = c(H\textsubscript{2}O) = 0.50 mol/L，c(CO\textsubscript{2}) = c(H\textsubscript{2}) = 0。
(a) 求平衡时各组分浓度。
(b) 若向平衡体系中再加入 0.30 mol CO\textsubscript{2}，求新平衡时 CO 的浓度。
(c) 讨论：增大 CO\textsubscript{2} 浓度后 CO 的转化率如何变化？这与勒夏特列原理是否一致？

\textbf{21.} \textbf{(全国初赛改编)} 已知 298 K 时：
- \(\mathrm{CaF_2(s) \rightleftharpoons Ca^{2+}(aq) + 2F^-(aq)}\)，\(K_{sp}^\theta = 3.9 \times 10^{-11}\)
- \(\mathrm{HF(aq) \rightleftharpoons H^+(aq) + F^-(aq)}\)，\(K_a^\theta = 6.8 \times 10^{-4}\)

计算 \(\mathrm{CaF_2(s) + 2H^+(aq) \rightleftharpoons Ca^{2+}(aq) + 2HF(aq)}\) 的 \(K^{\circ}\)。在 pH = 3 的缓冲溶液中，CaF\textsubscript{2} 的溶解度是多少？

\textbf{22.} \textbf{(综合计算)} 在 1000 K、100 kPa 下，反应 \(\mathrm{2HI(g) \rightleftharpoons H_2(g) + I_2(g)}\) 的 \(K^{\circ}\) = 0.016。若起始只有 1.00 mol HI 在 10 L 容器中：
(a) 求平衡时各组分的摩尔分数。
(b) 若在相同条件下起始有 1.00 mol HI + 0.50 mol H\textsubscript{2}，求 HI 的新转化率。
(c) 讨论增加 H\textsubscript{2} 对 HI 转化率的影响，与勒夏特列原理对照。

\textbf{23.} \textbf{(多平衡耦合)} 298 K 时已知：
- \(\mathrm{Ag_2CrO_4(s) \rightleftharpoons 2Ag^+(aq) + CrO_4^{2-}(aq)}\)，\(K_{sp}^\theta = 1.1 \times 10^{-12}\)
- \(\mathrm{Ag^+(aq) + 2NH_3(aq) \rightleftharpoons Ag(NH_3)_2^+(aq)}\)，\(K_f^\theta = 1.7 \times 10^7\)

\begin{enumerate}
\def\labelenumi{(\alph{enumi})}
\tightlist
\item
  计算 \(\mathrm{Ag_2CrO_4(s) + 4NH_3(aq) \rightleftharpoons 2Ag(NH_3)_2^+(aq) + CrO_4^{2-}(aq)}\) 的 \(K^{\circ}\)。
\item
  若向 1 L 2.0 mol/L 氨水中加入 0.01 mol Ag\textsubscript{2}CrO\textsubscript{4}，能否全部溶解？通过计算说明。
\item
  讨论此平衡在分析化学中''铬酸银溶于氨水''现象的意义。
\end{enumerate}

\textbf{24.} \textbf{(van't Hoff 综合)} 某反应在 300 K 时 \(K^{\circ}\) = 10，在 400 K 时 \(K^{\circ}\) = 100。
(a) 求该反应的 Δ\_rH°(假定为常数)。
(b) 求 Δ\_rS°(300 K 时)。
(c) 估算该反应的转变温度(Δ\_rG° = 0 时的温度)。
(d) 判断 500 K 时反应是否自发。

\textbf{25.} \textbf{(工业平衡综合)} 接触法制硫酸的关键反应：\(\mathrm{2SO_2(g) + O_2(g) \rightleftharpoons 2SO_3(g)}\)，Δ\_rH° = −198 kJ/mol。

{\def\LTcaptype{none} % do not increment counter
\begin{longtable}[]{@{}cc@{}}
\toprule\noalign{}
温度(K) & \(K^{\circ}\) \\
\midrule\noalign{}
\endhead
\bottomrule\noalign{}
\endlastfoot
700 & 2.4 × 10\textsuperscript{3} \\
800 & 1.9 × 10\textsuperscript{2} \\
900 & 3.2 × 10\textsuperscript{1} \\
\end{longtable}
}

\begin{enumerate}
\def\labelenumi{(\alph{enumi})}
\tightlist
\item
  计算 800 K 时 Δ\_rG°。
\item
  解释为什么工业上采用 400−500°C 而不是更低温度。
\item
  工业上采用''过量空气''(O\textsubscript{2} 过量)的目的是什么？从平衡转化率角度分析。
\item
  若初始气体组成为 SO\textsubscript{2}:O\textsubscript{2}:N\textsubscript{2} = 7:11:82(摩尔分数)，在 700 K、100 kPa 下求 SO\textsubscript{2} 的平衡转化率。
\end{enumerate}

\begin{center}\rule{0.5\linewidth}{0.5pt}\end{center}

\subsection{十一、练习题答案}\label{ux5341ux4e00ux7ec3ux4e60ux9898ux7b54ux6848}

\textbf{1.}
(a) \(K_c = \frac{[\mathrm{NH_3}]^2}{[\mathrm{N_2}][\mathrm{H_2}]^3}\)，\(K_p = \frac{p_{\mathrm{NH_3}}^2}{p_{\mathrm{N_2}} \cdot p_{\mathrm{H_2}}^3}\)
(b) \(K = p_{\mathrm{CO_2}}\)(纯固体不写入)
(c) \(K_c = \frac{[\mathrm{H^+}][\mathrm{CH_3COO^-}]}{[\mathrm{CH_3COOH}]}\)

\textbf{2.} \(K^{\circ}\) = e\^{}(-20000/(8.314×298)) = e\^{}(-8.07) = 3.1 × 10\textsuperscript{-4}。\(K^{\circ}\) ≪ 1，标准状态下正向反应几乎不发生。

\textbf{3.}
(a) 错误。K 只描述平衡位置，不描述速率。K 大不等于速率快。
(b) 错误。催化剂只改变速率，不改变 K。
(c) 正确。恒容充惰气，各组分分压不变，Q 不变，平衡不移动。
(d) 错误。放热反应升温 K 减小。

\textbf{4.} Δν = 2 − (2+1) = −1 \textless{} 0。增大总压 → Q \textless{} K → 平衡正向移动(向分子数减少方向)。

\textbf{5.} A(g) ⇌ 2B(g)，设转化 x mol/L：\(K_{\1}\) = (2x)\textsuperscript{2}/(1.0−x) = 4。解得 x = 0.33 mol/L。{[}A{]} = 0.67 mol/L，{[}B{]} = 0.67 mol/L。

\textbf{6.} 恒容充惰气：各组分分压不变(pᵢ = nᵢRT/V)，Q 不变，平衡不移动。恒压充惰气：体积增大，各组分分压降低(等效减压)，Q 改变，平衡向分子数增加方向移动。

\textbf{7.} K\textsubscript{3} = K\textsubscript{1}\textsuperscript{2} / K\textsubscript{2}

\textbf{8.} 减小。放热反应 Δ\_rH° \textless{} 0，由 van't Hoff 方程，升温 T\textsubscript{2} \textgreater{} T\textsubscript{1} 时 ln(K\textsubscript{2}/K\textsubscript{1}) = (Δ\_rH°/R)(1/T\textsubscript{1} − 1/T\textsubscript{2}) \textless{} 0，所以 K\textsubscript{2} \textless{} K\textsubscript{1}。

\textbf{9.} 加入 NO\textsubscript{2} 后 Q = {[}NO\textsubscript{2}{]}\textsuperscript{2}/{[}N\textsubscript{2}O\textsubscript{4}{]} 增大 → Q \textgreater{} K → 平衡逆向移动(向生成 N\textsubscript{2}O\textsubscript{4} 方向)。

\textbf{10.} \(K_{\1}\) = p\_B × p\_C。设分解压为 P，p\_B = p\_C = P/2。\(K_{\1}\) = (P/2)\textsuperscript{2} = P\textsuperscript{2}/4 = 0.04。P = √(0.16) = 0.40 kPa。

\begin{center}\rule{0.5\linewidth}{0.5pt}\end{center}

\textbf{11.} 设转化 x mol H\textsubscript{2}：\(K_{\1}\) = (2x/V)\textsuperscript{2}/{[}(1−x)/V × (1−x)/V{]} = 4x\textsuperscript{2}/(1−x)\textsuperscript{2} = 49。解得 x/(1−x) = 3.5，x = 0.778 mol。
平衡浓度：{[}H\textsubscript{2}{]} = {[}I\textsubscript{2}{]} = 0.0222 mol/L，{[}HI{]} = 0.156 mol/L。
α(H\textsubscript{2}) = 77.8\%。

\textbf{12.} Δν = 0，所以 \(K_{\1}\) = \(K_{\1}\) = 0.681。对于 Δν = 0 的反应，\(K_{\1}\) 与 \(K_{\1}\) 数值相等(因为 (RT)\^{}0 = 1)。

\textbf{13.} Δν = 0，\(K_{\1}\) = \(K^{\circ}\) = 4.08 × 10\textsuperscript{-4}。
NO 摩尔分数 = \(K^{\circ}\)/(1 + \(K^{\circ}\)) × \ldots{} 对于 Δν = 0，NO 的摩尔分数 = √\(K^{\circ}\)/(1 + √\(K^{\circ}\)) ≈ √\(K^{\circ}\)(因 \(K^{\circ}\) ≪ 1)。总压变化不影响 Δν = 0 的反应平衡组成。NO 摩尔分数不变。

\textbf{14.} \(K^{\circ}\) = \(K_{\1}\)° × \(K_{\1}\)° = 1.7×10\textsuperscript{7} × 1.8×10\textsuperscript{-10} = 3.1×10\textsuperscript{-3}。
\(K^{\circ}\) 较小但不极小。在浓氨水({[}NH\textsubscript{3}{]} ≈ 15 mol/L)中，Q ≪ K，反应可正向进行使 AgCl 溶解。在稀氨水中 {[}NH\textsubscript{3}{]} 不够大，Q \textgreater{} K，溶解不完全。

\textbf{15.} \(K^{\circ}\) = 1 时 \(\Delta \1^{\circ}\) = 0，即 \(\Delta \1^{\circ}\) = T\(\Delta \1^{\circ}\)。由 298 K 数据：\(\Delta \1^{\circ}\) = (\(\Delta \1^{\circ}\) − \(\Delta \1^{\circ}\))/T = (57200 − RT ln 0.148)/298 = (57200 + 4730)/298 = 207.6 J/(mol·K)。T = \(\Delta \1^{\circ}\)/\(\Delta \1^{\circ}\) = 57200/207.6 = 275 K ≈ 2°C。实际上 N\textsubscript{2}O\textsubscript{4}/NO\textsubscript{2} 体系在约 273 K 时 K ≈ 1。

\textbf{16.}
(a) Q = 0.50 × 0.50 = 0.25 \textgreater{} \(K_{\1}\) = 0.108，Q \textgreater{} K，不平衡。
(b) Q \textgreater{} K → 逆向自发 → NH\textsubscript{4}HS(s) 生成(气体结合沉积为固体)。

\textbf{17.}
(a) 设转化 x：\(K_{\1}\) = x\textsuperscript{2}/{[}(1−x)(1−x){]} = 4 → x = 0.667 mol/L。α(A) = 66.7\%。
(b) 设转化 x：\(K_{\1}\) = x\textsuperscript{2}/{[}(2−x)(0.5−x){]} = 4 → 3x\textsuperscript{2} − 10x + 4 = 0。x = 0.535。α(A) = 26.7\%。
增大一种反应物(A)的量会降低其自身的转化率，但提高另一种反应物(B)的转化率。这正是''过量一种反应物提高另一种转化率''的工业策略。

\textbf{18.}
(a) 低温 K 大(热力学有利)，但低温反应速率太慢(动力学不利)，生产效率极低。
(b) 铁催化剂降低活化能，提高反应速率；400−500°C 保证足够快的速率同时 K 不至于太小。
(c) 合成氨 Δν = 2−4 = −2 \textless{} 0，高压使平衡正向移动，提高 NH\textsubscript{3} 平衡浓度。

\textbf{19.}
(a) \(E^{\circ}\)(电池) = \(E^{\circ}\)(Fe\textsuperscript{3+}/Fe\textsuperscript{2+}) − \(E^{\circ}\)(S/H\textsubscript{2}S) = 0.77 − 0.14 = 0.63 V \textgreater{} 0，反应自发。
(b) \(K^{\circ}\) = e\^{}(nFE°/RT) = e\^{}(2×96485×0.63/(8.314×298)) = e\^{}49.2 ≈ 1.0 × 10\textsuperscript{21}
(c) \(K^{\circ}\) 极大，反应几乎完全。设平衡 {[}Fe\textsuperscript{3+}{]} = x：由电极电势计算 x ≈ 10\textsuperscript{-17} mol/L(极微量)。此时 Fe\textsubscript{2}S\textsubscript{3} 的离子积 Q = {[}Fe\textsuperscript{3+}{]}\textsuperscript{2}{[}S\textsuperscript{2-}{]}\textsuperscript{3} ≈ (10\textsuperscript{-17})\textsuperscript{2} × (0.05)\textsuperscript{3} ≈ 10\textsuperscript{-37} ≫ \(K_{\1}\) = 10\textsuperscript{-88}，说明 Fe\textsubscript{2}S\textsubscript{3} 沉淀会形成。

\begin{center}\rule{0.5\linewidth}{0.5pt}\end{center}

\textbf{20.}
(a) 设转化 x：\(K_{\1}\) = x\textsuperscript{2}/{[}(0.50−x)\textsuperscript{2}{]} = 5.0 → x = 0.50×√5/(1+√5) = 0.366 mol/L
{[}CO{]} = {[}H\textsubscript{2}O{]} = 0.134 mol/L，{[}CO\textsubscript{2}{]} = {[}H\textsubscript{2}{]} = 0.366 mol/L

\begin{enumerate}
\def\labelenumi{(\alph{enumi})}
\setcounter{enumi}{1}
\item
  新起始：{[}CO\textsubscript{2}{]} = 0.366+0.30 = 0.666，其余不变。设逆向反应 x'：\(K_{\1}\) = (0.366−x')\textsuperscript{2}/{[}(0.134+x')\textsuperscript{2}{]} = 5.0。解得 x' ≈ 0.042。新 {[}CO{]} ≈ 0.176 mol/L。
\item
  CO 转化率降低(从 73.2\% 降至约 64.8\%)。增大产物浓度使平衡逆向移动，CO 转化率下降。这与勒夏特列原理一致：增大 CO\textsubscript{2} 浓度(产物)，平衡向消耗 CO\textsubscript{2} 方向(逆向)移动。
\end{enumerate}

\textbf{21.} 目标反应 = CaF\textsubscript{2} 溶解 − 2×HF 电离的逆反应：
\(K^{\circ}\) = \(K_{\1}\)° / (\(K_{\1}\)°)\textsuperscript{2} = 3.9×10\textsuperscript{-11} / (6.8×10\textsuperscript{-4})\textsuperscript{2} = 8.4×10\textsuperscript{-5}

在 pH = 3 的缓冲溶液中 {[}H\textsuperscript{+}{]} = 10\textsuperscript{-3} mol/L。
设 CaF\textsubscript{2} 溶解度为 s：{[}Ca\textsuperscript{2+}{]} = s，{[}F\textsuperscript{-}{]} = 2s × (α\_F)，其中 α\_F = \(K_{\1}\)/(\(K_{\1}\)+{[}H\textsuperscript{+}{]}) ≈ 0.405
\(K_{\1}\)° = s × (2s × 0.405)\textsuperscript{2} = 0.656s\textsuperscript{3} = 3.9×10\textsuperscript{-11} → s ≈ 3.9×10\textsuperscript{-4} mol/L

\textbf{22.}
(a) 设转化 x：\(K_{\1}\) = \(K^{\circ}\) = 0.016(因 Δν = 0)。(x/(10))/(1−x)/10)\textsuperscript{2} = \ldots{} 简化：K = x\textsuperscript{2}/(1−x)\textsuperscript{2} = 0.016。x/(1−x) = 0.1265，x = 0.112。
摩尔分数：HI = 0.888/(0.888+0.056+0.056) = 88.8\%，H\textsubscript{2} = I\textsubscript{2} = 5.6\%。

\begin{enumerate}
\def\labelenumi{(\alph{enumi})}
\setcounter{enumi}{1}
\item
  新条件：起始 HI = 1, H\textsubscript{2} = 0.5, I\textsubscript{2} = 0。设 HI 变化 y：K = (0.5+y/10)(y/10)/{[}(1−y)/10{]}\textsuperscript{2} \ldots{} 较复杂。定性：H\textsubscript{2} 过量使平衡逆向移动，HI 转化率降低。
\item
  增加 H\textsubscript{2}(产物)→ Q \textgreater{} K → 平衡逆向移动 → HI 转化率降低。与勒夏特列原理一致。
\end{enumerate}

\textbf{23.}
(a) \(K^{\circ}\) = (\(K_{\1}\)°)\textsuperscript{2} × \(K_{\1}\)° = (1.7×10\textsuperscript{7})\textsuperscript{2} × 1.1×10\textsuperscript{-12} = 3.2×10\textsuperscript{2}

\begin{enumerate}
\def\labelenumi{(\alph{enumi})}
\setcounter{enumi}{1}
\item
  0.01 mol Ag\textsubscript{2}CrO\textsubscript{3} 在 1 L 2.0 mol/L NH\textsubscript{3} 中。反应消耗 0.04 mol NH\textsubscript{3}(极少量)，{[}NH\textsubscript{3}{]} ≈ 1.96 mol/L。
  Q = (0.02)\textsuperscript{2} × 0.01 / (1.96)\textsuperscript{4} = 4×10\textsuperscript{-6} / 14.8 = 2.7×10\textsuperscript{-7} ≪ \(K^{\circ}\) = 320
  Q ≪ K，反应正向进行极彻底，Ag\textsubscript{2}CrO\textsubscript{4} 能全部溶解。
\item
  银的铬酸盐溶于氨水是定性分析中鉴别 Ag\textsuperscript{+} 的经典反应。利用配位平衡(\(K_{\1}\) 大)驱动沉淀溶解平衡正向移动。
\end{enumerate}

\textbf{24.}
(a) ln(K\textsubscript{2}/K\textsubscript{1}) = ln(100/10) = 2.303 = (\(\Delta \1^{\circ}\)/R)(1/300 − 1/400) = \(\Delta \1^{\circ}\)/R × 8.33×10\textsuperscript{-4}
\(\Delta \1^{\circ}\) = 2.303 × 8.314 / 8.33×10\textsuperscript{-4} = 23,000 J/mol = 23.0 kJ/mol

\begin{enumerate}
\def\labelenumi{(\alph{enumi})}
\setcounter{enumi}{1}
\item
  \(\Delta \1^{\circ}\)(300 K) = −RT ln \(K^{\circ}\) = −8.314×300×2.303 = −5,744 J/mol = −5.74 kJ/mol
  \(\Delta \1^{\circ}\) = (\(\Delta \1^{\circ}\) − \(\Delta \1^{\circ}\))/T = (23000+5744)/300 = 95.8 J/(mol·K)
\item
  T\_转 = \(\Delta \1^{\circ}\)/\(\Delta \1^{\circ}\) = 23000/95.8 = 240 K
\item
  \(\Delta \1^{\circ}\)(500 K) = 23000 − 500×95.8 = 23000 − 47900 = −24,900 J/mol \textless{} 0，自发。
\end{enumerate}

\textbf{25.}
(a) Δ\_rG° = −RT ln \(K^{\circ}\) = −8.314×800×ln(1.9×10\textsuperscript{2}) = −8.314×800×5.25 = −34,900 J/mol ≈ −34.9 kJ/mol

\begin{enumerate}
\def\labelenumi{(\alph{enumi})}
\setcounter{enumi}{1}
\item
  低温 K 大(700 K 时 K = 2400)但速率慢；高温速率快但 K 小。400−500°C(约 700−800 K)是速率与平衡的最佳折中。
\item
  O\textsubscript{2} 过量使平衡正向移动(Q \textless{} K)，提高 SO\textsubscript{2} 的平衡转化率。过量空气成本低，是工业上提高转化率的经济手段。
\item
  设 SO\textsubscript{2} 转化 x mol：初始 7 mol SO\textsubscript{2}, 11 mol O\textsubscript{2}, 82 mol N\textsubscript{2}，共 100 mol。
  平衡：SO\textsubscript{2} = 7(1−x), O\textsubscript{2} = 11−3.5x, SO\textsubscript{3} = 7x, N\textsubscript{2} = 82。总 = 100−3.5x。
  \(p_i = y_i \times 100\) kPa。
\end{enumerate}

\(K^\theta = \frac{(7x/(100-3.5x))^2}{(7(1-x)/(100-3.5x))^2 \times ((11-3.5x)/(100-3.5x))} = \frac{x^2(100-3.5x)}{(1-x)^2(11-3.5x)} = 2400\)

迭代求解得 x ≈ 0.97(SO\textsubscript{2} 转化率约 97\%)。

\end{document}
